\section{ASPECTOS DE MERCADO}
\justifying

En la actualidad existen diversos dispositivos y laboratorios portátiles orientados a aplicaciones didácticas y de aprendizaje en electrónica de potencia. Un ejemplo es el desarrollo de K\&H PE-5000 Power Electronics Training System, \parencite{kandh-product}, que consiste en 28 módulos experimentales, un motor trifásico de jaula de ardilla, cargas diversas, control y dispositivos de medición. Dentro de los módulos incluye rectificadores monofásicos, trifásicos, inversores, entre otros. 

En la misma línea que el sistema anterior se encuentra el ofrecido por GW Instek, PTS-3100 \parencite{gwinstek-product}, que incorpora módulos como fuentes y cargas de DC y AC. Además, incluye un osciloscopio para poder realizar las mediciones de las experiencias. Otra solución destacada son los productos de la serie LabVolt de Festo Didactic, \parencite{festo-product}, los cuales abarcan módulos para aprender sobre teoría fundamental de corriente alterna, transistores de potencia y tiristores, comunicación analógica, entre otros.

Finalmente, se encuentran los módulos de entrenamiento ofrecidos por Eplimin S.A.C, \parencite{eplimin-product}, orientados a electrónica digital, electrónica de potencia y módulos analógicos, cada uno acompañado de un manual de usuario electrónico.

Cabe señalar que estas opciones comerciales no presentan precios en sus páginas oficiales, dado que se fabrican bajo pedido y están destinadas principalmente a instituciones educativas y centros de formación técnica.